\documentclass[11pt,a4paper]{article}
\usepackage[utf8]{inputenc}
\usepackage[T1]{fontenc}
\usepackage{amsfonts}
\usepackage[french]{babel}
\frenchbsetup{StandardLists=true} % à inclure si on utilise \usepackage[francais]{babel}
\usepackage{enumitem}
\usepackage{amssymb}
\usepackage{amsmath}
\usepackage[left=2cm,right=2cm,top=2cm,bottom=2cm]{geometry}
\usepackage{multirow}
\usepackage[dvipsnames]{xcolor}

\usepackage{fouriernc}
\usepackage{graphicx}

\usepackage{setspace}
\singlespacing
\setlength{\parindent}{0pt}

\usepackage{multicol}
\setlength{\columnseprule}{1pt}

\usepackage{tikz-osci}

\usepackage{fancyhdr}
\pagestyle{fancy}
\renewcommand\headrulewidth{1pt}
\fancyhead[L]{Lycée Denis Diderot }
\fancyhead[R]{Classe de Terminale}
\setcounter{page}{1}\fancyfoot[C]{page \thepage}
\fancyfoot[R]{Année 2025 2026}
\fancyfoot[L]{Alexis RUIZ}
\usepackage{multirow,tabularx,booktabs}
\newcommand{\cadreblanc}[1]{%
\medskip\framebox[\linewidth][c]{%
\rule{0mm}{#1mm}}\par
\medskip}

\renewcommand{\arraystretch}{2}
\newcommand*\acocher{\quitvmode{\fboxsep0pt \fboxrule0.8pt \fbox{\vrule height1.8ex depth.1ex width0pt \vrule height0pt depth0pt width1.9ex }}}

\begin{document}
%\hline 
\begin{center}
\vspace{-3mm}
\LARGE{Questionnaire : Tesla et la gueurre des courants} \\

\vspace{2mm}

\normalsize
\hline

\hrule\
\\[0,2cm]

\flushleft \textbf{NOM :\\[4mm]
Prénom:\\[0,2cm]
}
\hrule\
\\[0,2cm]


\normalsize
\end{center}

\vspace{2mm}

A partir de la BD de Sciences et Vie junior N°230 de novembre 2008, répondre aux questions ci-dessous : 


\begin{enumerate}
 
\item Quelle est la nationalité de Nicolas Tesla ?  

\cadreblanc{10}

\item Quel est son métier ?

\cadreblanc{10}

\item Décrire brièvement l'idée géniale de N Tesla ?

\cadreblanc{15}

\item Quels sont les inconvénients du courant continu ?

\cadreblanc{15}

\item  Quel personnage célèbre est parti rencontrer N Tesla aux États-Unis ?

\cadreblanc{10}

\item Quelle est la différence entre un rotor et un stator ?

\cadreblanc{20}

\item Quelle est la source du conflit entre N Tesla et T Edison qui finira par la démission du premier ?

\cadreblanc{15}

\item Quel est le nom de l'industriel qui va engager N Tesla comme consultant ?

\cadreblanc{10}

\item Quel est l'avantage de la tension alternative qui va lui permettre de s'imposer ?

\cadreblanc{20}

\item A quelle occasion la tension alternative va-t-elle définitivement s'imposer ?

\cadreblanc{10}


\item Dans quelle autre grande aventure scientifique Nicolas Tesla va-t-il alors se lancer ?

\cadreblanc{10}


\item Quel est le nom et la nationalité de son concurrent dans cette nouvelle aventure ?

\cadreblanc{10}


\item Combien de brevets industriels N tesla a-t-il déposé dans sa vie ?

\cadreblanc{15}


\end{enumerate}
\end{document}